\documentclass[12pt, a4paper]{scrreprt}
\usepackage[italian]{babel}

\usepackage{lecture}

\usepackage{acronyms}

\begin{document}

\makefront{}

\begingroup
    \hypersetup{linkbordercolor=white, linkcolor=black} 
    \tableofcontents

    \tcblistof[\section*]{def}{List of Definitions}
    \newpage
\endgroup

\chapter{Info}\label{chap:info} % (fold)
\begin{itemize}
    \item Matteo Golfarelli --- matteo.golfarelli@unibo.it
\end{itemize}
\paragraph{Esame}\label{par:esame} % (fold)
Prova orale in presenza sugli argomenti trattati a lezione.
Svolto a sportello su richiesta durante tutto l'anno in base alle seguenti regole:
\begin{itemize}
    \item comunicazione con circa 7-10 giorni di anticipo al docente la data a partire dalla quale avrà completato la preparazione;
    \item il docente comunica l'orario e la data in cui verrà sostenuto l'esame;
\end{itemize}
\paragraph{Obiettivi del Corso}\label{par:obiettivi_del_corso} % (fold)
\begin{itemize}
    \item conoscere il ruolo dei \ac{si} moderni e la loro evoluzione;
    \item conoscere le tipologie e la composizione dei moderni \ac{si} (portafoglio applicativo aziendale, ERP, CRM, Scada, IoT, Big Data);
    \item pianificare e innovare l'azienda tramite i sistemi informativi:
        \begin{itemize}
            \item analizzare e capire i processi aziendali;
            \item riprogettare i processi aziendali e i sistemi informatici;
            \item analizzare i rischi e costi/benefici della riprogettazione e degli investimenti in tecnologia;
            \item scrivere o rispondere a un capitolato tecnico;
        \end{itemize}
    \item condurre un progetto informatico nel settore \ac{si} (project management, gestione dei contratti);
\end{itemize}

% paragraph Obiettivi del corso (end)
% paragraph Esame (end)
% chapter Info (end)
%
% ======== INTRO ========
\import{sections/}{intro.tex}
%
\end{document}
